% Beta 9 manuscript section. This file is kept separate so it can be integrated
% into the venue-template version of main.tex without weakening the explicit
% artifact boundary.
\section{Machine-Checked Range-Analysis Soundness}

Magnitude-aware Change Capabilities depend on a quantitative analysis: if a
source change amount is summarized by an interval, concrete executions of the
supported expression must remain inside that interval. Beta 9 mechanizes this
obligation for a deliberately small integer fragment rather than treating all
production arithmetic as verified.

\subsection{Formal expression fragment}
The Lean module \texttt{PatchRange.lean} defines expressions
\[
e ::= n \mid x \mid e_1 + e_2 \mid e_1 - e_2 \mid -e \mid k \cdot e,
\]
where $n\in\mathbb{Z}$ and $k\in\mathbb{N}$. The final form is represented by
\texttt{scale k e}; restricting the factor to a non-negative integer constant
keeps this initial formal fragment linear while covering common contracts such
as \texttt{bonus * 2}.

An abstract environment $\Gamma$ maps variable names to ordered integer
intervals, while a concrete environment $\rho$ maps variable names to integer
values. The predicate \texttt{EnvRespects} states that each concrete value used
by $\rho$ lies inside the corresponding interval in $\Gamma$.

The formalization defines an executable abstract interpreter
\[
\mathit{analyzeRange}(e,\Gamma) : Option\ Interval
\]
and a concrete evaluator
\[
\mathit{evalRangeExpr}(e,\rho) : Option\ Int.
\]

\subsection{Soundness theorem}
Lean machine-checks the theorem \texttt{rangeAnalysisSound}:
\[
\begin{split}
&EnvRespects(\Gamma,\rho) \\
&\land\ analyzeRange(e,\Gamma)=Some(I) \\
&\land\ evalRangeExpr(e,\rho)=Some(v) \\
&\Longrightarrow v\in I.
\end{split}
\]
The proof is structural over the formal expression grammar. Addition,
subtraction and negation use the usual interval endpoint rules; constant scaling
is proved by induction through repeated addition. A concrete corollary proves
that $bonus\in[0,5]$ implies $2\cdot bonus\in[0,10]$.

\subsection{Production connection}
The normal JavaScript analyzer and the formal-certificate extractor are kept as
separate claim paths. \texttt{src/range-analysis.js} computes the ordinary
production interval. \texttt{src/formal-range.js} independently parses supported
production expression text into the formal \texttt{RangeExpr} vocabulary and
independently computes the corresponding interval. Certification requires these
intervals to agree.

For each accepted numeric change, the generated Lean certificate contains the
formal expression, its abstract environment and the claimed interval. Lean then
checks that \texttt{analyzeRange} returns exactly that interval and instantiates
\texttt{rangeAnalysisSound}. Formal CI exercises this path with a real Patch
program whose protected recipe is:

\begin{lstlisting}
allow reward:
  player.score may increase up to 10

make reward(player, bonus number 0..5):
  change player:
    add bonus * 2 to score
\end{lstlisting}

\subsection{Boundary}
This theorem does not establish soundness for every expression currently
accepted by the production analyzer. Beta 9 range certification intentionally
excludes division, decimal or floating-point semantics, and general
multiplication where neither operand is a non-negative integer literal. Such
constructs are not silently labeled verified.

Two important correspondence obligations remain. First, JavaScript parsing and
extraction from production source/AST into \texttt{RangeExpr} are independently
implemented and cross-checked but are not themselves machine-proved. Second,
the production runtime evaluator is not yet proved to correspond to
\texttt{evalRangeExpr}. The beta-9 result should therefore be stated as a
machine-checked soundness theorem for the formal integer fragment plus a
conservative production-validation boundary, not as full compiler or full
production-analyzer verification.
